\documentclass[a4paper, 11pt]{article}
\usepackage[utf8]{inputenc}
\usepackage[T1]{fontenc}
\usepackage[italian]{babel}
\usepackage{geometry}
\usepackage{graphicx}
\usepackage{tabularx}
\usepackage{fancyhdr}
\usepackage{hyperref}

% Margini
\geometry{top=3cm, bottom=3cm, left=2.5cm, right=2.5cm}

% Intestazione e piè di pagina
\pagestyle{fancy}
\fancyhf{}
\fancyhead[L]{\small \textit{0x06tudio} -- Verbale di Riunione}
\fancyhead[R]{\small Anno Accademico 2026/2027}
\fancyfoot[C]{\thepage}
\renewcommand{\headrulewidth}{0.4pt}

\begin{document}
\thispagestyle{plain}

\begin{minipage}{0.6\textwidth}
    \vspace{-2cm}
    {\huge \textbf{0x06tudio}} \\[0.5em]
    {\large Progetto didattico di Ingegneria del Software} \\[0.3em]
    {\normalsize Università degli Studi di Padova}
\end{minipage}
\begin{minipage}{0.4\textwidth}
    \raggedleft
    \vspace{0pt}
    \includegraphics[width=6cm]{../../media/logo.png}
    \vspace{1cm}
    \small
    \href{mailto:0x06tudio@gmail.com}{0x06tudio@gmail.com}
\end{minipage}

\vspace{0.5cm}

\begin{center}
    {\LARGE \textbf{Verbale Riunione}}
\end{center}

\vspace{0.5cm}


\begin{tabularx}{\textwidth}{l X}
    \textbf{Versione:} & 0.1.0 \\
    \textbf{Data:} & 2026/10/26 \\
    \textbf{Uso:} & Esterno \\
    \textbf{Partecipanti (Presenti):} &  Nicolò Beltrame, Pietro Brocadello, Alberto Cavalletto, Catalin Chirilov, Davide Colabove, Andrea Maggio\\
    \textbf{Destinatari:} & VarGroup Spa; Proff. Vardanega, Cardin \\
\end{tabularx}

\vspace{1cm}

\section*{Registro delle modifiche}
\begin{tabularx}{\textwidth}{|c|c|X|c|}
    \hline
    \textbf{Versione} & \textbf{Data} & \textbf{Descrizione} & \textbf{Autore} \\ \hline
    0.1.0 & 2026/10/26 & Creazione e stesura bozza iniziale & Andrea Maggio \\ \hline
\end{tabularx}

\newpage
\tableofcontents
\newpage


\section{Informazioni generali}
\begin{tabularx}{\textwidth}{l X}
    \textbf{Modalità / Luogo:} & Online / Microsoft Teams \\
    \textbf{Data svolgimento:} & 2026/10/26 \\
    \textbf{Ora inizio:} & 15:00 \\
    \textbf{Ora fine:} & 15:30 \\
    \textbf{Durata totale:} & 30 minuti \\
    \textbf{Scopo della riunione:} & Domande riguardo il capitolato C6 \\
    \textbf{Partecipanti esterni:} & Stefano Dindo, Laura Ferretti, Alessio Borgheresi, Michele Massaro (VarGroup SpA)
\end{tabularx}

\newpage


\section{Ordine del giorno}
Durante la riunione sono stati affrontati i seguenti punti:
\begin{enumerate}
    \item Come integrare il modello da sviluppare con liteLLM
    \item Valutare se dividere la cache tra team diversi
    \item Come effettuare il forecasting
    \item Come effettuare il caching di conversazioni composte da più messaggi

\end{enumerate}

\newpage


\section{Resoconto della discussione}

\subsection{Integrazione con liteLLM}
L'idea iniziale è effettuare un fork di liteLLM.
\subsection{Divisione della cache}
La cache è richiesta come univoca, perché in molti casi non è neanche ben definito chi appartiene a quale team.
Sarà tuttavia necessario dividere la dashboard poichè questa deve rendicontare come ogni sviluppatore spende le risorse che l'azienda mette a disposizione.
\subsection{Forecasting}
Si consiglia di usare dati storici per il forecasting, ma in fase esplorativa sarà possibile discutere alternative.
\subsection{Caching di conversazioni}
Rispetto alle singole richieste, fare il caching di una intera conversazione porterebbe ad un aumento della grandezza dei messaggi, e anche a maggiore possibilità di una cache miss anche se vengono effettuate domande simili. Consigliato valutare soluzioni di riduzione dinamica del contesto.
\newpage


\section{Approvazione Esterna}

Il presente verbale è stato condiviso con il referente aziendale in data YYYY/MM/DD ed è da ritenersi approvato e conforme a quanto discusso durante il colloquio. 

\vspace{3cm}
Firma del rappresentante:


\end{document}